\FloatBarrier
\section{The OECD cross-section}\label{app:oecd26}\setcounter{equation}{0}
This appendix documents the test of section \ref{sec:esc:crossSection}: the sample and the data, the U.S.\ at those data, where the political weight of the old and the cost parameter are calibrated, the second layer of the test, in which the political weight of the old is calibrated country by country, the test under other maintained assumptions, and the U.S., the UK and France as the deep test of section \ref{sec:esc:deepTest} and the cross-section measure them. Every economy is solved under log preferences with the cost of section \ref{sec:esc:cost}, whose parameter is $\lambda_{i,t} = \varepsilon\,\Pi_{i,t}/R_{i,t}$ at each date in the paper's specification, and exactly one candidate tax passes at every state and date of every solve (section \ref{sec:numerical}). The \oa{esc-oecd26} repeats the test on all 29 countries of figure \ref{fig:US:OECDdata}, with each country's own propensities to vote and with the U.S.\ $\lambda$ in every country.

\subsection{Sample and data}\label{app:oecd26:data}
The sample is the 29 countries of figure \ref{fig:US:OECDdata} less those in which more than a quarter of employment is informal in the ILO's measure (ILOSTAT, 2018--2020): Mexico at $56\%$, Turkey at $31\%$ and Korea at $29\%$. The next are the UK, the Czech Republic and the Netherlands, at $20\%$, $17\%$ and $14\%$, where the measure mostly reflects self-employment that their pension systems cover. Australia, Canada, Japan, New Zealand and the U.S.\ have no such measure and stay in. Korea and Turkey also stand apart in their pension systems: Korea's public pension dates only from 1988, and Turkey's retirement age for men is 52. On the 26 countries the relations in the data keep the signs they have on the 29, except the flat relation between spending and the Gini (\oa{esc-oecd26-all}).

Every country, the U.S.\ included, is measured in the same way. The number of workers per retiree $\nu_t$ is the thirty-year growth of the total population in the medium variant of the United Nations' World Population Prospects 2024, from 1990 to 2080, with 1960 at the value of 1990 and the periods after 2080 at the value of 2080, as in the calibrations of section \ref{sec:oecd}; for France it is the plain ratio, without the correction for the retirement age of appendix \ref{app:US}. Pension spending and the index $I$ are those of figure \ref{fig:US:OECDdata}, and the observed design is $\theta = 2-1/I$, the inverse of the mapping in section \ref{sec:esc:crossSection}, clipped to $[0,1]$. Australia, whose means-tested pension puts its index just below one half, and Ireland are therefore at the Beveridgean corner, and Austria, Finland, France, Hungary, Italy and Spain at the Bismarckian one. The income distribution is that of the gross earnings of men in full-time employment in the OECD earnings database, in 2019 or the nearest year, through a lognormal fitted separately below and above the median to the ratios of the first and ninth deciles to the median. Households are grouped at the U.S.\ income percentiles of section \ref{sec:oecd}, so the three groups hold $58\%$, $18\%$ and $24\%$ of households in every country. Everything else takes the U.S.\ values of section \ref{sec:oecd}: $\alpha = \xi = 0.30$, the propensities to vote from the CPS, a taste for leisure common to all groups, which under log preferences sets hours but no rate, and the discount factor of the calibration without the cost, $\beta = 0.655$, imposed as on the UK and France in appendix \ref{app:US}. Table \ref{table:OECD26:omega} lists the inputs that differ across countries.

\subsection{The U.S.\ at these data}\label{app:oecd26:anchor}
At these inputs the U.S.\ differs from its calibration of section \ref{sec:oecd} in two respects. Its design is $\theta = 0.735$, the inverse of its index, $39.2/49.6 = 0.790$, where section \ref{sec:esc} has $0.728$, since the means of the CPS groups lie close to, but not exactly at, half the mean and the mean. Its dispersion of relative incomes is $0.333$, against $0.436$ at the CPS groups: the OECD series covers men in full-time work only and is fitted to three points of the distribution, so it misses the part-time and part-year earners that the CPS records. With $\beta$ imposed, we calibrate $\omega$ to the U.S.\ tax rate and $\varepsilon$ so that the electorate of 1990 chooses the observed design for 2020, as in section \ref{sec:esc:calibration}, which gives $\omega = 1.294$ and $\varepsilon = 20.014$, against $20.445$ at the CPS groups (table \ref{table:US_ESC:country}). In the first layer of the test every country takes these values (figure \ref{fig:OECD26:layer1}).

\subsection{Each country's political weight}\label{app:oecd26:layer2}
In the second layer each country keeps the U.S.\ $\beta$ and $\varepsilon$, and its political weight of the old is calibrated so that its tax rate of 2020 at the observed design matches its pension spending, as $\omega$ is in section \ref{sec:oecd}. Its design is then predicted as in the first layer, which extends the test of section \ref{sec:esc:deepTest} from the UK and France to 25 countries. Table \ref{table:OECD26:omega} reports the political weights and the predicted designs. The weight of the old runs from $0.947$ in Hungary to $1.758$ in Luxembourg, against $1.294$ in the U.S. The predicted designs run from $0.618$ in Ireland to $0.924$ in Finland, a narrower range than the observed one, and they keep the relations of the first layer with population growth and spending, while their relation with the Gini is no longer significant (table \ref{table:OECD26:slopes}). For 17 of the 25 countries the predicted design lies on the same side of the U.S.\ design as the observed one. The UK is among the exceptions, predicted at $0.791$, above the U.S.\ at $0.735$ (section \ref{sec:esc:deepTest}). France is predicted at $0.879$, short of its corner but about half of the way to it from the U.S.\ design, where the deep test closes little of that distance: the two measurements differ in France's workers per retiree, its political weight of the old and its propensities to vote (table \ref{table:OECD26:reconcile}).

%% GENERATED by python/paper/build.py from results/ -- do not edit by hand.
%% Rebuild:  .venv\Scripts\python.exe python\paper\build.py --only OECD26_Omega
%% Source:   results/oecd26/layer2.csv, anchor.csv; data/oecdFigure1.csv, oecd26_sample.csv
\begin{table}[!htb]
\centering
\begin{threeparttable}
\caption{The countries of the cross-country test, each calibrated to its pension spending}
\label{table:OECD26:omega}
\renewcommand{\arraystretch}{1.25}
\footnotesize
\begin{tabularx}{\textwidth}{lYYYYYYYY}
\toprule
\textbf{Country} & $\nu_{2020}$ & $\tilde V$ & \textbf{Spending} & $I$ & $\theta$ & $\omega$ & $\theta$ \textbf{chosen} & $I$ \textbf{chosen} \\
\midrule 
Australia & 1.50 & 0.202 & 4.2 & 0.499 & 0.000 & 1.331 & 0.801 & 0.834 \\ % row: aus
Austria & 1.16 & 0.161 & 14.3 & 1.000 & 1.000 & 1.610 & 0.922 & 0.928 \\ % row: aut
Belgium & 1.16 & 0.104 & 11.5 & 0.643 & 0.445 & 1.459 & 0.932 & 0.936 \\ % row: bel
Canada & 1.37 & 0.171 & 6.4 & 0.729 & 0.629 & 1.291 & 0.826 & 0.852 \\ % row: can
Czech Republic & 1.02 & 0.170 & 8.6 & 0.603 & 0.343 & 1.141 & 0.872 & 0.887 \\ % row: cze
Denmark & 1.13 & 0.116 & 8.2 & 0.639 & 0.436 & 1.203 & 0.878 & 0.891 \\ % row: dnk
Finland & 1.11 & 0.094 & 12.6 & 1.000 & 1.000 & 1.411 & 0.949 & 0.952 \\ % row: fin
France & 1.16 & 0.150 & 14.4 & 1.000 & 1.000 & 1.610 & 0.923 & 0.928 \\ % row: fra
Germany & 1.05 & 0.148 & 11.2 & 0.892 & 0.880 & 1.247 & 0.898 & 0.907 \\ % row: deu
Greece & 1.04 & 0.225 & 17.9 & 0.857 & 0.833 & 1.709 & 0.924 & 0.930 \\ % row: grc
Hungary & 0.94 & 0.265 & 7.7 & 1.000 & 1.000 & 0.947 & 0.797 & 0.831 \\ % row: hun
Iceland & 1.44 & 0.112 & 3.0 & 0.711 & 0.593 & 1.145 & 0.823 & 0.849 \\ % row: isl
Ireland & 1.42 & 0.240 & 3.3 & 0.500 & 0.000 & 1.155 & 0.748 & 0.799 \\ % row: irl
Italy & 1.05 & 0.134 & 17.4 & 1.000 & 1.000 & 1.688 & 0.951 & 0.953 \\ % row: ita
Japan & 1.02 & 0.124 & 9.6 & 0.750 & 0.667 & 1.136 & 0.902 & 0.911 \\ % row: jpn
Luxembourg & 1.65 & 0.187 & 9.0 & 0.847 & 0.820 & 1.761 & 0.876 & 0.890 \\ % row: lux
Netherlands & 1.17 & 0.116 & 6.8 & 0.953 & 0.951 & 1.129 & 0.852 & 0.871 \\ % row: nld
New Zealand & 1.49 & 0.141 & 5.1 & 0.604 & 0.344 & 1.345 & 0.856 & 0.874 \\ % row: nzl
Norway & 1.27 & 0.101 & 7.6 & 0.759 & 0.683 & 1.274 & 0.898 & 0.908 \\ % row: nor
Poland & 1.00 & 0.240 & 11.4 & 0.962 & 0.961 & 1.203 & 0.873 & 0.888 \\ % row: pol
Portugal & 1.04 & 0.249 & 13.6 & 0.982 & 0.981 & 1.386 & 0.876 & 0.890 \\ % row: prt
Spain & 1.22 & 0.183 & 12.8 & 1.000 & 1.000 & 1.571 & 0.903 & 0.912 \\ % row: esp
Sweden & 1.21 & 0.081 & 8.5 & 0.868 & 0.848 & 1.266 & 0.925 & 0.930 \\ % row: swe
Switzerland & 1.29 & 0.127 & 6.9 & 0.831 & 0.796 & 1.240 & 0.872 & 0.886 \\ % row: che
United Kingdom & 1.17 & 0.185 & 7.8 & 0.694 & 0.559 & 1.199 & 0.823 & 0.850 \\ % row: gbr
United States$^{\ast}$ & 1.34 & 0.333 & 7.4 & 0.790 & 0.735 & 1.295 & 0.805 & 0.837 \\ % row: usa
\bottomrule
\end{tabularx}
\begin{tablenotes}[flushleft]
\footnotesize
\item[] \textit{Note:} $\nu_{2020}$, workers per retiree; $\tilde V$, the dispersion of relative incomes at the U.S.\ income percentiles; spending in \% of GDP; $I$, the Bismarckian index, and $\theta = 2 - 1/I$ clipped to $[0,1]$. $\omega$ is calibrated to spending at $\beta = 0.655$ and the U.S.\ $\lambda = 11.218$; the last two columns are the design chosen in 2020 and its index. $^{\ast}$The U.S.\ is the anchor, where $\lambda$ is calibrated.
\end{tablenotes}
\end{threeparttable}
\end{table}


\subsection{Other maintained assumptions}
Table \ref{table:OECD26:robustness} repeats the relations of the design under each combination of the two transports of the cost, $\lambda$ and $\varepsilon$, two rules for $\beta$ and two sources of the propensities to vote. Under the second rule each country's $\beta$ is set so that its return in 2020 is the U.S.\ one, with $\beta$ and $\omega$ calibrated together in the second layer. With each country's own propensities to vote, from the CSES, the U.S.\ anchor is recalibrated at its own CSES profile, and Luxembourg, with no CSES election, leaves data and model alike. The index on the Gini is where the propensities to vote decide. At the U.S.\ ones the index falls with the Gini in the first layer under either cost and either rule, from $-0.67$ to $-0.78$, and significantly, and in the second layer by less and not significantly. At each country's own it is about zero in the first layer, from $-0.04$ to $0.09$, and positive in the second, from $0.40$ to $0.62$, the sign of the data, though not significantly. In the paper's specification at each country's own propensities, spending also rises with the index at about four times the data's slope, and in the first layer the index falls with population growth at about a third of the data's slope on those 25 countries, $-0.22$ against $-0.62$. The rule for $\beta$ moves no slope of the index in the first layer by more than $0.05$. Under $\lambda$ the index falls more steeply with population growth than under $\varepsilon$ in every comparison, and its relations with the Gini keep their significance or its absence. In every row the index falls with population growth and spending rises with the index, as in the data, and significantly, except for the fall of the index in the second layer under $\varepsilon$ with a common return, $-0.10$.

%% GENERATED by python/paper/build.py from results/ -- do not edit by hand.
%% Rebuild:  .venv\Scripts\python.exe python\paper\build.py --only OECD26_Robustness
%% Source:   results/oecd26/layer1.csv, layer2.csv, anchor.csv; data/oecdFigure1.csv, oecd26_sample.csv
\begin{table}[!htb]
\centering
\begin{threeparttable}
\caption{The design relations under other maintained assumptions}
\label{table:OECD26:robustness}
\renewcommand{\arraystretch}{1.25}
\footnotesize\setlength{\tabcolsep}{3pt}
\begin{tabularx}{\textwidth}{lYYYYYY}
\toprule
 & \multicolumn{2}{c}{\textbf{Index on growth}} & \multicolumn{2}{c}{\textbf{Index on the Gini}} & \multicolumn{2}{c}{\textbf{Spending on the index}} \\
\cmidrule(lr){2-3}\cmidrule(lr){4-5}\cmidrule(lr){6-7}
\textbf{Cost, rule} & Layer 1 & Layer 2 & Layer 1 & Layer 2 & Layer 1 & Layer 2 \\
\midrule 
\multicolumn{7}{l}{\textit{U.S.\ turnout, 26 countries}} \\
Data & \multicolumn{2}{c}{$-0.45$ $(0.16)$} & \multicolumn{2}{c}{$0.89$ $(0.98)$} & \multicolumn{2}{c}{$16.33$ $(3.67)$} \\ % row: cps_data
$\lambda$, common $\beta$ & $-0.28$ $(0.04)$ & $-0.19$ $(0.06)$ & $-0.73$ $(0.33)$ & $-0.19$ $(0.38)$ & $39.84$ $(6.25)$ & $51.27$ $(7.36)$ \\ % row: cps_size_common
$\lambda$, common $R$ & $-0.31$ $(0.04)$ & $-0.19$ $(0.06)$ & $-0.78$ $(0.37)$ & $-0.19$ $(0.39)$ & $39.23$ $(6.13)$ & $50.89$ $(7.39)$ \\ % row: cps_size_commonR
$\varepsilon$, common $\beta$$^{\ast}$ & $-0.23$ $(0.04)$ & $-0.15$ $(0.05)$ & $-0.67$ $(0.29)$ & $-0.24$ $(0.31)$ & $42.06$ $(8.02)$ & $58.14$ $(10.35)$ \\ % row: cps_sizePiR_common
$\varepsilon$, common $R$ & $-0.23$ $(0.04)$ & $-0.10$ $(0.06)$ & $-0.67$ $(0.30)$ & $-0.04$ $(0.36)$ & $41.85$ $(8.76)$ & $50.47$ $(8.83)$ \\ % row: cps_sizePiR_commonR
\multicolumn{7}{l}{\textit{Own turnout, 25 countries}} \\
Data & \multicolumn{2}{c}{$-0.62$ $(0.18)$} & \multicolumn{2}{c}{$0.87$ $(1.01)$} & \multicolumn{2}{c}{$16.39$ $(3.75)$} \\ % row: cses_data
$\lambda$, common $\beta$ & $-0.29$ $(0.02)$ & $-0.35$ $(0.07)$ & $-0.02$ $(0.30)$ & $0.42$ $(0.46)$ & $51.41$ $(4.00)$ & $48.55$ $(4.55)$ \\ % row: cses_size_common
$\lambda$, common $R$ & $-0.30$ $(0.02)$ & $-0.35$ $(0.07)$ & $-0.04$ $(0.31)$ & $0.42$ $(0.46)$ & $52.68$ $(4.18)$ & $48.52$ $(4.53)$ \\ % row: cses_size_commonR
$\varepsilon$, common $\beta$ & $-0.22$ $(0.02)$ & $-0.27$ $(0.05)$ & $0.06$ $(0.23)$ & $0.40$ $(0.35)$ & $66.59$ $(5.45)$ & $61.94$ $(6.94)$ \\ % row: cses_sizePiR_common
$\varepsilon$, common $R$ & $-0.20$ $(0.02)$ & $-0.23$ $(0.07)$ & $0.09$ $(0.21)$ & $0.62$ $(0.38)$ & $76.18$ $(7.18)$ & $56.20$ $(5.29)$ \\ % row: cses_sizePiR_commonR
\bottomrule
\end{tabularx}
\begin{tablenotes}[flushleft]
\footnotesize
\item[] \textit{Note:} Least-squares slopes with standard errors in parentheses; the regressors are those of figure \ref{fig:US:OECDdata}, the model\textquotesingle s index is $1/(2-\theta)$ of the design in force in 2020 on the freely simulated path, chosen in 1990, and its spending is $\tau(1-\alpha)$ at the observed design. Rows: the cost parameter carried from the U.S.\ anchor to every country, its $\lambda$ or its $\varepsilon$ with $\lambda_i = \varepsilon\,\Pi_{i,2020}/R_{i,2020}$, and $\beta$, the U.S.\ value in every country (common $\beta$) or set so that each country\textquotesingle s 2020 return is the U.S.\ one (common $R$, with $\beta$ and $\omega$ calibrated together in layer 2). Own turnout: every country\textquotesingle s propensities to vote from the CSES, the U.S.\ anchor recalibrated at its own CSES profile (Luxembourg, with no CSES election, left out). $^{\ast}$The paper\textquotesingle s specification. --: not computed.
\end{tablenotes}
\end{threeparttable}
\end{table}


\subsection{The U.S., the UK and France measured twice}
Table \ref{table:OECD26:reconcile} sets the three economies of the deep test, each calibrated from its own household data, beside the same countries in the cross-section, measured from the OECD series and grouped at the U.S.\ income percentiles. The U.S.\ design is $0.728$ at the CPS groups and $0.735$ at the OECD data (appendix \ref{app:oecd26:anchor}). The dispersion of relative incomes is lower in the OECD series for all three countries, $0.333$ against $0.436$ for the U.S., $0.185$ against $0.223$ for the UK and $0.150$ against $0.218$ for France. The cross-section gives every country the U.S.\ propensities to vote, where the deep test gives the UK and France their own. And France has $1.04$ workers per retiree in 2020 in the deep test, with the correction for its earlier retirement age (appendix \ref{app:US}), against $1.16$ in the cross-section, which takes the plain ratio for every country. This last difference alone separates France's two political weights of the old, $1.446$ and $1.610$ under every cost. At France's Bismarckian corner the cost vanishes and every retiree's pension and savings are proportional to its earnings, so under log preferences every type supports the tax alike, and the tax depends neither on the income distribution nor on the propensities to vote (appendices \ref{app:PEE:LOG} and \ref{app:esc:cost}); with more workers per retiree, the same spending then takes a larger weight of the old (proposition \ref{prop:LOG:PEE}). Under the U.S.\ $\varepsilon$ the two measurements put France's design at $0.743$ and $0.879$ and the UK's at $0.673$ and $0.791$. At each country's own propensities to vote the cross-section brings both most of the way to the deep test's designs (\oa{esc-oecd26-turnout}).

%% GENERATED by python/paper/build.py from results/ -- do not edit by hand.
%% Rebuild:  .venv\Scripts\python.exe python\paper\build.py --only OECD26_Reconcile
%% Source:   results/paper/usCalibrationSummary.csv; results/esc/escCalibration{,Costs}.csv, escCountry{,Costs}.csv; results/oecd26/anchor.csv, layer2.csv
\begin{table}[!htb]
\centering
\begin{threeparttable}
\caption{The U.S., the UK and France in the deep test and in the cross-section}
\label{table:OECD26:reconcile}
\renewcommand{\arraystretch}{1.25}
\begin{tabularx}{\textwidth}{lYYYYYY}
\toprule
 & \multicolumn{3}{c}{\textbf{Deep test}} & \multicolumn{3}{c}{\textbf{Cross-section}} \\
\cmidrule(lr){2-4}\cmidrule(lr){5-7}
 & \textbf{U.S.} & \textbf{UK} & \textbf{France} & \textbf{U.S.} & \textbf{UK} & \textbf{France} \\
\midrule 
$\theta$ observed & 0.728 & 0.560 & 1.000 & 0.735 & 0.559 & 1.000 \\ % row: theta
$\tilde V$ & 0.436 & 0.223 & 0.218 & 0.333 & 0.185 & 0.150 \\ % row: vtilde
$\nu_{2020}$ & 1.34 & 1.17 & 1.04 & 1.34 & 1.17 & 1.16 \\ % row: nu
$\omega$, no cost & 1.231 & 1.103 & 1.446 & 1.259 & 1.145 & 1.610 \\ % row: omega_none
$\omega$, U.S.\ $\lambda$ & 1.280 & 1.166 & 1.446 & 1.294 & 1.198 & 1.610 \\ % row: omega_size
$\theta$ chosen, U.S.\ $\lambda$ & 0.728 & 0.690 & 0.795 & 0.735 & 0.805 & 0.903 \\ % row: chosen_size
$\omega$, U.S.\ $\varepsilon$ & 1.280 & 1.163 & 1.446 & 1.294 & 1.196 & 1.610 \\ % row: omega_sizePiR
$\theta$ chosen, U.S.\ $\varepsilon$ & 0.728 & 0.673 & 0.743 & 0.735 & 0.791 & 0.879 \\ % row: chosen_sizePiR
\bottomrule
\end{tabularx}
\begin{tablenotes}[flushleft]
\footnotesize
\item[] \textit{Note:} Deep test: each country\textquotesingle s own calibration, the U.S.\ in section \ref{sec:oecd} and the UK and France in appendix \ref{app:US} ($\theta$, $\nu_{2020}$ and $\omega$ without the cost), and $\tilde V$, $\omega$ and the design chosen at the U.S.\ $\lambda$ and $\varepsilon$ from the test of table \ref{table:US_ESC:country}; for the U.S.\ from the calibration of its cost, which re-elects its observed design. Cross-section: the OECD data of appendix \ref{app:oecd26}, the U.S.\ anchor and the second layer, $\omega$ calibrated to each country\textquotesingle s spending at the U.S.\ $\beta$ and cost parameter. The design chosen is the one in force in 2020, chosen in 1990, on the freely simulated path; under $\varepsilon$, $\lambda_i = \varepsilon\,\Pi_{i,2020}/R_{i,2020}$. --: not computed.
\end{tablenotes}
\end{threeparttable}
\end{table}

